\begin{afterword}
\summaryhead

The post opens with Bill Hibbard's 2001 proposal: train simple machines to recognize happiness in human faces, voices and body language, then hard-wire the result as the values of more intelligent machines. The author had asked whether such an AI might tile the galaxy with molecular smiley faces; Hibbard replied that a superintelligent recognizer would not be fooled by distinctions trivial to humans.

The author answers with a parable, told with the warning that it may not be true: an army neural network that learned to tell cloudy days from sunny ones instead of tanks from forest. The number of concepts compatible with any finite set of labeled examples is enormous, so which one a learner adopts depends on its inductive bias, and which one is correct depends on the user's values. New borderline cases, like a molecular smiley face or Terri Schiavo, raise questions the training data did not test.

Hibbard's reply, the author says, underestimates the complexity of value-laden concepts and assumes that a superintelligence would share Hibbard's ranking. The author calls this the fallacy of magical categories. Friendly AI is a problem of communicating category boundaries, and even a perfect smile recognizer would lead a superintelligence to manufacture smiles.

\responsehead

The central argument is sound, and it answers Hibbard's reply on its own terms. Hibbard said a superintelligence could make ``discriminations that are trivial to current humans.'' The post separates telling two images apart from classifying a new image under a concept. A learner can see that a molecular smiley face differs from a human face and still count both as smiles, because nothing in the training data says which way to generalize. That choice is made by the learner's inductive bias, the assumptions it brings to new cases. The counting behind it is correct, and it is the textbook argument the book already gave in ``Superexponential Conceptspace.'' The lesson, that success on training data guarantees nothing once the real cases come from a different process, is the standard assumption of supervised learning, stated clearly.

The parable is weaker than the argument, and the post says so. Gwern Branwen's search for the tank story found it to be an urban legend, probably based on a question asked in the 1960s. The post had flagged that it ``might or might not be fact,'' and nothing in the argument depends on it. Hibbard had objected in 2006 that pairing Hibbard's proposal with the army's networks conflated two different things. The post pairs them again, but its main argument applies to any learner, so the objection does not affect it.

One smaller point. The post says Hibbard's 2001 piece was ``published in a peer-reviewed journal''; Hibbard's own page calls it a column, and I could not confirm that it was peer reviewed, though the book that followed shows the position is not a strawman.

In short: a sound argument that training data do not fix how a learner classifies new cases, which answers Hibbard's reply directly, introduced by a tank parable that, as the post allows, is probably a legend.
\end{afterword}
